\writeSectionFrame{\presentationTitle}{Das Brücken-Muster}
\begin{frame}{Steckbrief}
	\begin{itemize}
 		\item Deutscher Name: Brücke
 		\item Englischer Name: Bridge
 		\item Gruppe: Strukturmuster
 	\end{itemize}
\end{frame}

\begin{frame}{Brücke}
	\begin{block}{Zweck}
 		Entkoppeln einer Abstraktion von ihrer Implementierung, so dass beide unabhängig voneinander variiert werden können.
 	\end{block}
\end{frame}


\frameWithImageOnly{\BILDERPFAD/bridge1.jpg}
\note{
	An dieser Stelle kommt das Brückenmuster ins Spiel. Es schlägt eine Brücke
	zwischen einem Interface und einer Implementierung, wobei ein Interface nicht
	im Java-Sinne gemeint es. Es kann sich auch um eine abstrakte Klasse handeln.
	Wir bauen also nicht unterschiedliche Klassen für unsere Komponenten, sondern
	wir bauen eine zweite Hierarchie mit Implementierern auf.  
}

\begin{frame}{Die Brücke löst das Problem}
	\begin{center}
 		\myImageWithSource{0.45}{\BILDERPFAD/bridge2.png}{https://refactoring.guru/design-patterns/bridge}
 	\end{center}
\end{frame}

\note{ 
 	Mit dem Brückenmuster können wir die Abstraktion und die Implementierung getrennt voneinander
 	ändern. Und die Brücke verbindet sie miteinander. 
}

\begin{frame}{Allgemeines Klassendiagramm}
    \begin{tikzpicture}[scale=0.6, transform shape]
        \umlclass[type=abstract, x=6, y=0]{Client}{
        }{
        }
    
        \umlclass[type=abstract, x=2, y=0]{Abstraction}{
        }{
            +operation() : void
        }
        
        % Verfeinerte Abstraktionen
        \umlclass[x=6, y=4]{RefinedAbstractionA}{}{
            +operation() : void
        }
        \umlclass[x=6, y=-4]{RefinedAbstractionB}{}{
            +operation() : void
        }
        
        % Implementor Schnittstelle
        \umlclass[type=interface, x=-4, y=0]{Implementor}{
            +operationImpl() : void
        }
        
        % Konkrete Implementierungen
        \umlclass[x=-8, y=4]{ConcreteImplementorA}{}{
            +operationImpl() : void
        }
        \umlclass[x=-8, y=-4]{ConcreteImplementorB}{}{
            +operationImpl() : void
        }
        
        % Verbindungen
        \umluniassoc{Client}{Abstraction}
        \umlaggreg{Abstraction}{Implementor} % Komposition
        \umlinherit{RefinedAbstractionA}{Abstraction}
        \umlinherit{RefinedAbstractionB}{Abstraction}
        \umlinherit{ConcreteImplementorA}{Implementor}
        \umlinherit{ConcreteImplementorB}{Implementor}
    \end{tikzpicture}
\end{frame}

\note{
	Die Abstraktion definiert die Schnittstelle der Abstraktion und enthält eine Referenz auf die 
	Implementierer-Klasse. Sie wird erweitert durch die KonkreteAbstraktion, die neues Verhalten
	ins Spiel bringen können. Die Verbindung zwischen Abstraktion und Implementierer kann zur
	Laufzeit geschehen und ist deswegen als Aggregation modelliert. Sie ist eine eher lose
	Kopplung. Die Implementierer-Hierarchie muss nicht 1:1 der Schnittstellen-Abstraktion
	entsprechen, denn beide sollen sich ja unabhängig voneinander entwickeln. Es kann also
	auch weitere Implementiererklassen geben.
}

\begin{frame}{Beispiel ohne Brücke}
    \begin{tikzpicture}[scale=0.6, transform shape]
        % Interfaces
        \umlinterface[x=0, y=6]{Shape}{}{}
        \umlinterface[x=0, y=3]{Color}{}{}
        
        % Klassen
        \umlclass[x=-3, y=-2]{RedColorTriangle}{}{}
        \umlclass[x=3, y=-2]{RedColorPentagon}{}{}
        
        \umlclass[x=-6, y=0]{GreenColorTriangle}{}{}
        \umlclass[x=6, y=0]{GreenColorPentagon}{}{}
        
        \umlclass[x=-4, y=3]{Triangle}{}{}
        \umlclass[x=4, y=3]{Pentagon}{}{}
        
        % Verbindungen
        \umlinherit{Triangle}{Shape}
        \umlinherit{Pentagon}{Shape}
        
        \umlinherit{RedColorTriangle}{Triangle}
        \umlinherit{GreenColorTriangle}{Triangle}
        
        \umlinherit{RedColorPentagon}{Pentagon}
        \umlinherit{GreenColorPentagon}{Pentagon}
        
        \umlinherit{RedColorTriangle}{Color}
        \umlinherit{GreenColorTriangle}{Color}
        
        \umlinherit{RedColorPentagon}{Color}
        \umlinherit{GreenColorPentagon}{Color}
    \end{tikzpicture}
\end{frame}

\begin{frame}{Klassendiagramm mit Brücke}
    \begin{tikzpicture}[scale=0.6, transform shape]
        \umlclass[x=-9, y=0]{Client}{
        }{
        }
    
        \umlclass[type=abstract, x=2, y=0]{Color}{
        }{
            +applyColor(): void
        }
        
        % Verfeinerte Abstraktionen
        \umlclass[x=6, y=4]{RedColor}{}{
            +applyColor(): void
        }
        \umlclass[x=0, y=-4]{GreenColor}{}{
            +applyColor(): void
        }

        \umlclass[x=6, y=-4]{BlueColor}{}{
            +applyColor(): void
        }
        
        % Implementor Schnittstelle
        \umlclass[type=interface, x=-4, y=0]{Shape}{
            +Shape(color: Color): void \\
            +applyColor(): void
        }
        
        % Konkrete Implementierungen
        \umlclass[x=-8, y=-4]{Circle}{}{
            +Circle(color: Color): void \\
            +applyColor(): void
        }
        \umlclass[x=-8, y=4]{Rectangle}{}{
            +Rectangle(color: Color): void \\
            +applyColor(): void
        }
        
        % Verbindungen
        \umluniassoc{Client}{Shape}
        \umluniassoc{Shape}{Color}
        \umlinherit{GreenColor}{Color}
        \umlinherit{RedColor}{Color}
        \umlinherit{BlueColor}{Color}
        \umlinherit{Rectangle}{Shape}
        \umlinherit{Circle}{Shape}
    \end{tikzpicture}
\end{frame}

\begin{frame}[fragile]{Oberklasse Implementierung}
	\begin{exampleblock}{shapes.Color}
 		\begin{lstlisting}
public interface Color {
	void applyColor();
}
 		\end{lstlisting}
 	\end{exampleblock}
\end{frame}

\begin{frame}[fragile]{Konkrete Implementierung}
	\begin{exampleblock}{RedColor}
 		\begin{lstlisting}
public class RedColor implements Color {
	@Override
	public void applyColor() {
		System.out.println("Rote Farbe wird verwendet");
	}
	
	@Override
	public String toString() {
		return "rot";
	}
}
 		\end{lstlisting}
 	\end{exampleblock}
\end{frame}

\begin{frame}[fragile]{Oberklasse Abstraktion}
	\begin{exampleblock}{Shape}
 		\begin{lstlisting}
public abstract class Shape {
	private Color color;
	
	public Shape(Color color) {
		this.color = color;
	}
	
	public void changeColor(Color newColor) {
		this.color = newColor;
	}
	
	public abstract void applyColor();
	
	public Color getColor() {
		return color;
	}
}
 		\end{lstlisting}
 	\end{exampleblock}
\end{frame}

\begin{frame}[fragile]{Konkrete Abstraktion}
	\begin{exampleblock}{Circle}
 		\begin{lstlisting}
public class Circle extends Shape {

	public Circle(Color color) {
		super(color);
	}

	@Override
	public void applyColor() {
		System.out.println("Unser Kreis hat die Farbe: "
            + getColor());		
	}
}
 		\end{lstlisting}
 	\end{exampleblock}
\end{frame}

\frameWithImageOnly{\BILDERPFAD/bridge3.jpg}
\begin{frame}{Klassendiagramm -- Fernbedienung}
    \begin{tikzpicture}[scale=0.6, transform shape]
        % Klassen
        \umlclass[x=0, y=0]{RemoteControlUser}{}{}
        
        \umlclass[x=6, y=0]{RemoteControl}{
            - Device device
        }{
            + toggleOn() : void \\
            + toggleOff() : void
        }
        
        \umlclass[x=12, y=0]{Device}{
            - on : boolean \\
            - volume : int
        }{
            + turnOn() : void \\
            + turnOff() : void \\
            + increaseVolume() : void \\
            + decreaseVolume() : void
        }
        
        \umlclass[x=6, y=-3]{AdvancedRemoteControl}{}{
            + mute() : void
        }
        
        \umlclass[x=10, y=-6]{Radio}{}{}
        
        \umlclass[x=14, y=-6]{TV}{}{}
        
        % Verbindungen
        \umluniassoc[arg=, pos=0.5]{RemoteControlUser}{RemoteControl} % Unidirektionale Assoziation
        \umlaggreg[arg=, pos=0.5]{RemoteControl}{Device} % Aggregation (weißer Diamant)
        \umlinherit{AdvancedRemoteControl}{RemoteControl} % Vererbung
        \umlinherit{Radio}{Device} % Vererbung
        \umlinherit{TV}{Device} % Vererbung
        
        \end{tikzpicture}    
\end{frame}

\begin{frame}[fragile]{Oberklasse Implementierung}
	\begin{exampleblock}{Device}
 		\begin{lstlisting}
public abstract class Device {
	private boolean on;
	private int volume;
	
	public void turnOn() {
		on = true;
	}
 
	public void increaseVolume() {
		volume++;
	}
	
	public void decreaseVolume() {
		if (volume > 0) {
			volume--;
		}
	}
}
 		\end{lstlisting}
 	\end{exampleblock}
\end{frame}

\begin{frame}[fragile]{Konkrete Implementierung}
	\begin{exampleblock}{Radio}
 		\begin{lstlisting}
public class Radio extends Device {
	private int radioStation;
	
	public void nextRadioStation() {
		radioStation++;
	}
	
	public void prevRadioStation() {
		if (radioStation > 0) {
			radioStation--;
		}
	}
	
	public int getRadioStation() {
		return radioStation;
	}
}
 		\end{lstlisting}
 	\end{exampleblock}
\end{frame}

\begin{frame}[fragile]{Oberklasse Abstraktion}
	\begin{exampleblock}{RemoteControl}
 		\begin{lstlisting}
public class RemoteControl {
	private Device currentDevice;
	private List<Device> devices;

	public void switchDevice() {
		if (currentIndex < devices.size() - 1) {
			currentIndex++;
		} else {
			currentIndex = 0;
		}
		currentDevice = devices.get(currentIndex);
	}
	
	public void toggleOn() {
		currentDevice.turnOn();
	}
}
 		\end{lstlisting}
 	\end{exampleblock}
\end{frame}

\begin{frame}[fragile]{Konkrete Abstraktion}
	\begin{exampleblock}{AdvancedRemoteControl}
 		\begin{lstlisting}
public class AdvancedRemoteControl extends RemoteControl {
	public AdvancedRemoteControl() {
		super();
	}
	
	public void mute() {
		Device device = getDevice(); 
		int volume = device.getVolume();
		
		for (int i = volume; i > 0; i--) {
			device.decreaseVolume();
		}
		
		System.out.println("...");
	}
}

 		\end{lstlisting}
 	\end{exampleblock}
\end{frame}

\begin{frame}{Klassendiagramm -- Zeichen-APIs}
    \begin{tikzpicture}[scale=0.5, transform shape]
        \umlclass[x=-6, y=14]{Client}{}{
        }
    
        \umlinterface[x=0, y=10]{DrawingAPI}{
            + drawCircle(x, y, radius) : void \\
            + drawRectangle(x, y, width, height) : void
        }{}
        
        \umlclass[x=-3, y=6]{DrawingAPI1}{}{
            + drawCircle(x, y, radius) : void \\
            + drawRectangle(x, y, width, height) : void
        }
        
        \umlclass[x=6, y=4]{DrawingAPI2}{}{
            + drawCircle(x, y, radius) : void \\
            + drawRectangle(x, y, width, height) : void
        }
        
        % Abstraktionsebene
        \umlclass[x=0, y=14]{Shape}{
            - DrawingAPI drawingAPI
        }{
            + draw() : void \\
            + resize(factor : float) : void
        }
        
        \umlclass[x=10, y=8]{Circle}{
            - x : float \\
            - y : float \\
            - radius : float
        }{
            + draw() : void \\
            + resize(factor : float) : void
        }
        
        \umlclass[x=10, y=12]{Rectangle}{
            - x : float \\
            - y : float \\
            - width : float \\
            - height : float
        }{
            + draw() : void \\
            + resize(factor : float) : void
        }
        
        % Verbindungen
        \umlinherit{DrawingAPI1}{DrawingAPI} % Vererbung
        \umlinherit{DrawingAPI2}{DrawingAPI} % Vererbung
        
        \umlaggreg[arg=, pos=0.5]{Shape}{DrawingAPI} % Komposition
        
        \umlinherit{Circle}{Shape} % Vererbung
        \umlinherit{Rectangle}{Shape} % Vererbung
        \umluniassoc{Client}{Shape}
        
    \end{tikzpicture}    
\end{frame}

\begin{frame}[fragile]{Oberklasse Abstraktion}
	\begin{exampleblock}{APIShape}
 		\begin{lstlisting}
public abstract class APIShape {
    protected DrawingAPI drawingAPI;

    protected APIShape(DrawingAPI drawingAPI) {
        this.drawingAPI = drawingAPI;
    }

    public abstract void draw();
}
 		\end{lstlisting}
 	\end{exampleblock}
\end{frame}

\begin{frame}[fragile]{Konkrete Abstraktion 1}
	\begin{exampleblock}{APICircle}
 		\begin{lstlisting}
public class APICircle extends APIShape {
    private double x, y, radius;

    public APICircle(double x, double y, 
            double radius, DrawingAPI drawingAPI) {
        super(drawingAPI);
        this.x = x;
        this.y = y;
        this.radius = radius;
    }

    @Override
    public void draw() {
        drawingAPI.drawCircle(x, y, radius);
    }
}
 		\end{lstlisting}
 	\end{exampleblock}
\end{frame}

\begin{frame}[fragile]{Konkrete Abstraktion 2}
	\begin{exampleblock}{APIRectangle}
 		\begin{lstlisting}
public class APIRectangle extends APIShape {
    private double x, y, width, height;

    public APIRectangle(double x, double y, double width,
            double height, DrawingAPI drawingAPI) {
        super(drawingAPI);
        this.x = x;
        this.y = y;
        this.width = width;
        this.height = height;
    }

    @Override
    public void draw() {
        drawingAPI.drawRectangle(x, y, width, height);
    }
}
 		\end{lstlisting}
 	\end{exampleblock}
\end{frame}

\begin{frame}[fragile]{Oberklasse Implementierung}
	\begin{exampleblock}{DrawingAPI}
 		\begin{lstlisting}
public interface DrawingAPI {
    void drawCircle(double x, double y, double radius);
    void drawRectangle(double x, double y, 
        double width, double height);
}
 		\end{lstlisting}
 	\end{exampleblock}
\end{frame}

\begin{frame}[fragile]{Konkrete Implementierung 1}
	\begin{exampleblock}{OpenGLAPI}
 		\begin{lstlisting}
public class OpenGLAPI implements DrawingAPI {
    @Override
    public void drawCircle(double x, double y, 
            double radius) {
        System.out.println("...");
    }

    @Override
    public void drawRectangle(double x, double y, double width, 
            double height) {
        System.out.println("...");
    }
}
 		\end{lstlisting}
 	\end{exampleblock}
\end{frame}

\begin{frame}[fragile]{Konkrete Implementierung 2}
	\begin{exampleblock}{DirectXAPI}
 		\begin{lstlisting}
public class DirectXAPI implements DrawingAPI {
    @Override
    public void drawCircle(double x, double y, double radius) {
        System.out.println("...");
    }

    @Override
    public void drawRectangle(double x, double y, double width, 
            double height) {
        System.out.println("...");
    }
}
 		\end{lstlisting}
 	\end{exampleblock}
\end{frame}

\begin{frame}{Klassendiagramm -- Benachrichtigungssystem}
    \begin{tikzpicture}[scale=0.5, transform shape]
        \umlclass[x=-6, y=14]{Client}{}{
        }
    
        \umlclass[type=abstract, x=0, y=10]{Message}{
            \#Message(MessageSender messageSender) \\
            \textit{+void send(content: String)}
        }{}
        
        \umlclass[x=-3, y=6]{TextMessage}{}{
            \#TextMessage(MessageSender messageSender) \\
            +void send(content: String)
        }
        
        \umlclass[x=6, y=4]{AlertMessage}{}{
            \#TextMessage(MessageSender messageSender) \\
            +void send(content: String)
        }
        
        % Abstraktionsebene
        \umlinterface[x=0, y=14]{MessageSender}{
        }{
            +sendMessage(message: String): void;
        }
        
        \umlclass[x=10, y=8]{SMSSender}{
        }{
            +sendMessage(message: String): void;
        }
        
        \umlclass[x=10, y=14]{EmailSender}{
        }{
            +sendMessage(message: String): void;
        }

        \umlclass[x=10, y=12]{PushNotificationSender}{
        }{
            +sendMessage(message: String): void;
        }        
        
        % Verbindungen
        \umlinherit{PushNotificationSender}{MessageSender} % Vererbung
        \umlinherit{EmailSender}{MessageSender} % Vererbung
        \umlinherit{SMSSender}{MessageSender} % Vererbung
        
        \umluniassoc[arg=, pos=0.5]{Message}{MessageSender} % Komposition
        \umluniassoc[arg=, pos=0.5]{Client}{Message} % Komposition
        
        \umlinherit{TextMessage}{Message} % Vererbung
        \umlinherit{AlertMessage}{Message} % Vererbung
        
    \end{tikzpicture}    
\end{frame}

\begin{frame}[fragile]{Abstraktion}
	\begin{exampleblock}{Message}
 		\begin{lstlisting}
public abstract class Message {
    private MessageSender messageSender;

    protected Message(MessageSender messageSender) {
        this.messageSender = messageSender;
    }
    
    public abstract void send(String content);
    
    protected MessageSender getMessageSender() {
		return messageSender;
	}
}

 		\end{lstlisting}
 	\end{exampleblock}
\end{frame}

\begin{frame}[fragile]{Konkrete Abstraktion}
	\begin{exampleblock}{Message}
 		\begin{lstlisting}
public class AlertMessage extends Message {
    public AlertMessage(MessageSender messageSender) {
        super(messageSender);
    }

    @Override
    public void send(String content) {
        getMessageSender().sendMessage("...");
    }
}
 		\end{lstlisting}
 	\end{exampleblock}
\end{frame}

\begin{frame}[fragile]{Oberklasse Implementierung}
	\begin{exampleblock}{MessageSender}
 		\begin{lstlisting}
public interface MessageSender {
    void sendMessage(String message);
}
 		\end{lstlisting}
 	\end{exampleblock}
\end{frame}

\begin{frame}[fragile]{Oberklasse Implementierung}
	\begin{exampleblock}{MessageSender}
 		\begin{lstlisting}
public class SMSSender implements MessageSender {
    @Override
    public void sendMessage(String message) {
        System.out.println("...");
    }
}
 		\end{lstlisting}
 	\end{exampleblock}
\end{frame}

\begin{frame}[fragile]{Klient}
	\begin{exampleblock}{MonitoringSystem}
 		\begin{lstlisting}
public class MonitoringSystem {
	private List<MessageSender> activeSenders;
	private List<MessageSender> inactiveSenders;
	
	public void sendTextMessage() {
		for (MessageSender messageSender: activeSenders) {
			messageSender.sendMessage("Eine Textnachricht");
		}
	}
	
	public void activateSender(int index) {
		if (index < inactiveSenders.size()) {
			// ...
		}
	}
 		\end{lstlisting}
 	\end{exampleblock}
\end{frame}